% -*- coding: utf-8 -*-
%UTF-8: äöüß
% !TEX root = ./main.tex
%preface.tex

\usepackage[utf8]{inputenc}
\usepackage[T1]{fontenc}
\usepackage[english,ngerman]{babel}
\usepackage{color}
\usepackage{graphicx}
\usepackage{rotating}
\usepackage{xspace} 
\usepackage{epstopdf}
\usepackage{hyperref}
\usepackage{tikz}
\usetikzlibrary{positioning}
\usetikzlibrary{shapes}

\usepackage{multicol}
\usepackage{algorithm}
\usepackage{algpseudocode}
\renewcommand{\algorithmicrequire}{\textbf{Input:}}
\renewcommand{\algorithmicensure}{\textbf{Output:}}
\floatname{algorithm}{Algorithmus}
% declaration of the new block
\algblock{ParDo}{EndParDo}
% customising the new block
\algnewcommand\algorithmicparfor{\textbf{for}}
\algnewcommand\algorithmicpardo{\textbf{pardo}}
\algnewcommand\algorithmicendparfor{\textbf{end\ pardo}}
\algrenewtext{ParDo}[1]{\algorithmicparfor\ #1\ \algorithmicpardo}
\algrenewtext{EndParDo}{\algorithmicendparfor}

\newcommand{\Renew}{\textsc{Renew}}


%Eigene Befehle
%Alles was zwischen den Klammern kommt, das wird nicht angezeigt:
%\commentout{DAS HIER WIRD NICHT SICHTBAR}
\newcommand{\commentout}[1]{}

\newcommand{\longpage}{\enlargethispage{\baselineskip}}
\newcommand{\Longpage}{\enlargethispage{2\baselineskip}}
\newcommand{\code}[1]{\texttt{#1}}
\newcommand{\source}[1]{\code{#1}}

\newcommand\plainEPS[2][0.5]{%
  \centerline{%
    \includegraphics[scale=#1]{#2.eps}%
    }%
  }%

\newcommand\mydoubleEPS[3][0.5]{
  \begin{figure}[htbp]
    \begin{minipage}{0,5\textwidth}
      \centerline{
        \includegraphics[scale=#1]{#2.eps}
      }
    \end{minipage}%
    \begin{minipage}{0,5\textwidth}
      \centerline{
        \includegraphics[scale=#1]{#3.eps}
      }
    \end{minipage}
  \end{figure}%
}     



%%%Einige Einstellungen
\mode<presentation>
{
  \usetheme{Malmoe}
%  \usecolortheme[named=DarkGreen]{structure}
  \usecolortheme{seagull}
  \useinnertheme{default}

% Disable navigation bar
  \setbeamertemplate{navigation symbols}{}

%\setbeamercovered{transparent}
% oder auch nicht
}

%Beispiel für \commentout{}
%Alternatives Layout
\commentout{
\mode<presentation>
{
  \usetheme{Warsaw}
  \usecolortheme{seahorse}
  \useinnertheme{circles}

	\setbeamercovered{
		still covered={\opaqueness<1>{15}\opaqueness<2>{10}\opaqueness<3>{5}\opaqueness<4->{2}},
		again covered={\opaqueness<1->{15}}
	}

  \setbeamertemplate{navigation symbols}{\insertframenavigationsymbol}
}
}%Ende alternative Layout



% Der Pfeil ist fuer Aufzaehlungen gedacht, die eine Folgerung
% aus vorigen Punkten sind.
\newcommand{\pfeil}{$\rhd$}

% Logo
%\pgfdeclareimage[height=0.5cm]{uni-logo}{images/UHH_Logo_only}
%\logo{\pgfuseimage{uni-logo}}
\logo{\includegraphics[height=0.5cm]{content/images/UHH_Logo_only.pdf}}
